% Continuous model-theoretic chapter. Shared macros supplied by the main paper.
\providecommand{\val}[1]{\lVert #1\rVert}
\providecommand{\qval}{\operatorname{qval}}
\providecommand{\Amp}{\operatorname{Amp}}
\providecommand{\CM}{\mathcal M}
\providecommand{\CN}{\mathcal N}
\section{Continuous model theory from finite correspondences}
\label{sec:continuous-models}

The repository already contains presented categories, a presentation-wise
Yoneda theorem, finite limits and colimits, and a setoid interpretation of
intuitionistic first-order derivations. Repeating those constructions would
not solve the present analytic problem. The contribution of this section is
a quantitative interpretation of \emph{real-valued formulas on continuous
structures}. Its central estimate is unusually economical: a two-sided
correspondence with atomic error $\eta$ preserves every formula built from
$\min,\max,1-(-),\inf,\sup$ with the \emph{same} error $\eta$.
Quantifier alternation does not multiply the error.

This is a restricted, effective part of continuous model theory. In the usual
semantic setting, predicates take values in $[0,1]$ and quantifiers are
infima and suprema; uniform continuity makes passage to completions
meaningful. These conventions and their general model-theoretic setting
are described in Ben Yaacov--Berenstein--Henson--Usvyatsov,
\emph{Model theory for metric structures}, Sections 2--4
(\url{https://math.univ-lyon1.fr/~begnac/articles/mtfms.pdf}).
Here we supply finite algorithms and arithmetic error proofs for specified
compact presentations. We do not import that work's compactness,
ultraproduct, saturation, or model-existence theorems into iPRA.

\subsection{The syntax and the scope of a presentation}

Fix a finite collection of sorts and a finite collection of atomic
recipes. An atomic recipe $P$ has a specified finite list of argument sorts.
For a semantic interpretation its value is a uniformly continuous function
into $[0,1]$. A formula is obtained from atomic recipes and rational constants
in $[0,1]$ by the operations
\begin{equation}
 \min(\varphi,\psi),\qquad
 \max(\varphi,\psi),\qquad
 1-\varphi,\qquad
 \inf_{x:s}\varphi,\qquad
 \sup_{x:s}\varphi .
 \label{eq:cm-syntax}
\end{equation}
Free variables carry sorts. Quantification always uses a nonempty sort.
The operations in \eqref{eq:cm-syntax} are numerical operations on formula
values, not replacements for the intuitionistic implication and
quantifiers of the ambient arithmetic.

The fixed-atom convention includes ordinary terms. For a fixed formula,
an occurrence $P(t_1(\bar x),\ldots,t_r(\bar x))$ can first be compiled into
one atomic recipe, with its own approximation procedure and modulus.
The syntax tree of each $t_i$ is fixed in advance. This avoids requiring a
finite grid to be closed under the structure's operations. Such a grid
usually is not a substructure: rounded addition on a grid, for example,
need not satisfy the original algebraic identities.
The approximation evaluates the original composite predicates at grid
points. It does not replace the operations by arbitrary rounded tables.

For clarity we first use the basis \eqref{eq:cm-syntax}.
One may also use a specified rational-evaluable connective
$C:[0,1]^r\to[0,1]$ with a proved rational Lipschitz bound $L_C$ for the
maximum metric. Define its formula amplification by
\begin{equation}
\begin{split}
 A_P&=1,\qquad A_c=0,\\
 A_{\min(\varphi,\psi)}=A_{\max(\varphi,\psi)}
   &=\max(A_\varphi,A_\psi),\qquad A_{1-\varphi}=A_\varphi,\\
 A_{\inf_x\varphi}=A_{\sup_x\varphi}&=A_\varphi,\\
 A_{C(\varphi_1,\ldots,\varphi_r)}
   &=L_C\max_i A_{\varphi_i}.
\end{split}
\label{eq:cm-amplification}
\end{equation}
Use $A=\max(1,A_\varphi)$ when choosing a precision budget.
For the core basis, $A_\varphi\leq1$ regardless of formula size or
quantifier depth. Truncated addition is permitted by taking $L_C=2$;
averaging has $L_C=1$. A separate proof or modulus is needed for any
additional connective. Arbitrary continuous connectives are not silently
treated as primitive recursive.

\subsection{The sharp finite transfer theorem}

Let $A_s$ and $B_s$ be nonempty finite sets for each sort $s$. Atomic
interpretations are rational-valued tables. For finite structures,
$\inf$ and $\sup$ mean the minimum and maximum of a finite nonempty list.
A \emph{total correspondence} consists of relations
$R_s\subseteq A_s\times B_s$ such that every member of either carrier
has a related member of the other. Assume
\begin{equation}
 \bigl|P^A(a_1,\ldots,a_r)-P^B(b_1,\ldots,b_r)\bigr|
 \leq\eta
 \quad\text{whenever }a_iR_{s_i}b_i\text{ for every }i.
 \label{eq:cm-atomic-distortion}
\end{equation}
The hypothesis concerns \emph{all related tuples}, not merely tuples
chosen by one matching algorithm. This distinction matters when a variable
occurs more than once or several variables are quantified.

\begin{theorem}[Finite approximate elementarity]
\label{thm:cm-transfer}
For every fixed formula $\varphi(\bar x)$ and coordinatewise related
assignments $\bar a,\bar b$,
\begin{equation}
 \bigl|\varphi^A(\bar a)-\varphi^B(\bar b)\bigr|
 \leq A_\varphi\eta.
 \label{eq:cm-transfer}
\end{equation}
In particular the bound is $\eta$ for the basis
\eqref{eq:cm-syntax}. The finite correspondence and all numerical tests
can be coded by natural numbers. For each fixed formula this theorem
has an iPRA proof using rational arithmetic, bounded scans, and
quantifier-free induction.
\end{theorem}

\begin{proof}
For two pairs of numbers satisfying $|u-u'|,|v-v'|\leq e$, the inequalities
$u\leq u'+e$ and $v\leq v'+e$ imply
$\max(u,v)\leq\max(u',v')+e$. Interchanging the pairs proves the
reverse inequality. The same argument, or complementation, proves the
minimum case. Also
$|(1-u)-(1-u')|=|u-u'|$. A connective $C$ obeys the asserted estimate by
its Lipschitz hypothesis.

It remains to prove the quantifier step without introducing an extra
error. Write $u(a)$ and $v(b)$ for the two subformula values after fixing
the already related parameters. Suppose $|u(a)-v(b)|\leq e$ whenever
$aRb$. Choose an index $a_*$ attaining $\max_Au$ by a finite scan.
Totality gives $b$ with $a_*Rb$, and therefore
\[
 \max_Au=u(a_*)\leq v(b)+e\leq\max_Bv+e.
\]
Choose an index attaining $\max_Bv$ and use the other half of totality
to obtain the reverse inequality. Thus the maxima differ by at most
$e$. For minima, choose $a_*$ attaining $\min_Au$ and a related $b$:
$\min_Bv\leq v(b)\leq u(a_*)+e=\min_Au+e$.
The reverse inequality follows from a minimizer on $B$.
This proves both quantifier cases with no increase in $e$.
Structural induction on the fixed syntax proves
\eqref{eq:cm-transfer}.

For the arithmetic implementation, represent a finite set by its length
and entries, a relation by a Boolean matrix, and rational values by
signed numerator and positive denominator. For the core basis a common
denominator reduces the proof to inequalities on natural numerators.
A finite scan stores its current extremal value and the first index
attaining it. Its invariant says that the index is in the scanned prefix,
its value equals the current extremum, and every scanned entry lies on
the appropriate side of that value. All these assertions are bounded.
After Booleanization they form a quantifier-free induction formula.
The uses of totality eliminate bounded existential witnesses; they do
not choose an infinite family.
Formula induction is a metatheoretic construction of one iPRA proof per
fixed syntax tree. No induction on an unbounded semantic satisfaction
predicate is asserted.
\end{proof}

Two-sided totality is essential. A one-point subset of $[0,1]$ preserves
the distance predicate on its own points exactly, but its value of
$\sup_x\sup_y d(x,y)$ is zero, whereas the interval's value is one.
The missing part of the comparison is coverage of the interval.

\begin{proposition}[Composition of approximate interpretations]
\label{prop:cm-compose}
If $R$ is a total correspondence of atomic error $\eta$ from $A$ to $B$
and $S$ one of error $\theta$ from $B$ to $C$, their relational composite
is total and has atomic error at most $\eta+\theta$.
\end{proposition}
\begin{proof}
For $a$ choose a related $b$ and then a related $c$; the reverse direction
uses the other two totality assertions. For a tuple related by the
composite, take the finitely many intermediate entries $b_i$ witnessing
its coordinates. Applying the two atomic inequalities and adding them
gives the stated bound. Finite relational composition is a bounded
disjunction over the middle carrier, so its table is primitive recursive.
Associativity follows by reassociating the two finite existential
quantifiers; the diagonal relation is an identity. Reversing a
correspondence preserves its error.
\end{proof}

Consequently finite presentations and total correspondences form a
category whose errors compose additively. The value of each sentence
is stable under its arrows with the amplification in
\eqref{eq:cm-amplification}. This supplies an abstract comparison
calculus for continuous observations. It does not assert that every
correspondence is the graph of a continuous map, or that arbitrary
metric spaces have been assembled into an internal category of iPRA.

\subsection{Primitive-recursive compact presentations}
\label{subsec:cm-presentations}

For each sort $s$, take a rational pseudometric $d_s$ on numerical
codes $D_s$ with diameter at most one, and a fixed primitive-recursive
procedure producing a nonempty finite list $N_s(n)$ of codes. If a
validity predicate is needed, all operations are totalized outside it,
and every assertion below is guarded by validity. The central net
contract is
\begin{equation}
 \forall a\in D_s\ \exists j<|N_s(n)|\quad
 d_s(a,N_s(n)_j)\leq 2^{-n}.
 \label{eq:cm-net}
\end{equation}
Its proof is supplied in iPRA; it is not inferred by inspecting a
finite list. Because distances are rational, a least suitable index
can alternatively be found by a bounded scan.

For each fixed atomic recipe $P$, supply rational approximants
$p_P(\bar a,p)\in[0,1]$ and a nondecreasing primitive-recursive
modulus $\mu_P$. Their arithmetical contracts are coherence and stability:
\begin{equation}
\begin{split}
 |p_P(\bar a,p)-p_P(\bar a,q)|
   &\leq 2^{-p}+2^{-q},\\
 d_\infty(\bar a,\bar b)\leq2^{-\mu_P(h)}
 \quad&\Longrightarrow\quad
 |p_P(\bar a,p)-p_P(\bar b,q)|
   \leq2^{-p}+2^{-q}+2^{-h}.
\end{split}
\label{eq:cm-atom-contract}
\end{equation}
Here $d_\infty$ is the maximum of the coordinate distances, and the
maximum over an empty tuple is zero. These are first-order statements
about fixed numerical functions. Denominators are cleared before
placing them in the object language. One can require dyadic
approximants, with a fixed extra factor for the formula's rational
constants, to make the common-denominator implementation immediate.

The usual semantic description follows constructively: the coherent
sequence $p_P(\bar a,p)$ is a real number; its value is uniformly
continuous on the dense presentation and hence extends to the
completion. This interpretation takes place outside first-order iPRA.
It is not an enumeration of all real numbers by primitive-recursive
program codes.

Let
\begin{equation}
 U(h)=\max\bigl(h,\mu_P(j):
       P\text{ occurs in the fixed formula},\ j\leq h\bigr).
 \label{eq:cm-common-modulus}
\end{equation}
This is a primitive-recursive, nondecreasing common modulus. The
additional maximum over $j\leq h$ makes monotonicity explicit even if
the original moduli were not normalized.

Given dense free parameters $\bar a$, append their entries to the
appropriate lists $N_s(n)$, obtaining finite nonempty carriers
$K_s(n,\bar a)$. Duplicates cause no problem. Interpret every atom
on these carriers by $p_P(-,p)$ and evaluate the formula by finite
minimum and maximum. Denote the resulting rational number by
\begin{equation}
 E_\varphi(\bar a;n,p).
 \label{eq:cm-evaluator}
\end{equation}
The same carrier is used at every quantifier in a single evaluation.
Appending free parameters permits the evaluation to be described
literally as a finite structure with an assignment. Finite supports
may contain any additional specified points: none of the estimates
depends on which extra points were appended.

\begin{theorem}[Cross-mesh estimate]
\label{thm:cm-cross-mesh}
If $n,m\geq U(h)$, then
\begin{equation}
 |E_\varphi(\bar a;n,p)-E_\varphi(\bar a;m,q)|
 \leq A_\varphi\bigl(2^{-p}+2^{-q}+2^{-h}\bigr).
 \label{eq:cm-cross-mesh}
\end{equation}
More generally the same estimate compares parameters $\bar a,\bar b$
when $d_\infty(\bar a,\bar b)\leq2^{-U(h)}$.
\end{theorem}
\begin{proof}
Relate an entry of the first carrier to an entry of the second whenever
their distance is at most $2^{-\min(n,m)}$, in the equal-parameter
case. Every entry has a partner by \eqref{eq:cm-net}: the net contained
in the other carrier covers \emph{all} dense codes, including any
appended parameters. Thus the relation is total both ways.
Identical free parameters are related.
The inequality $n,m\geq U(h)$ and
\eqref{eq:cm-atom-contract} give atomic distortion at most
$2^{-p}+2^{-q}+2^{-h}$ on every related tuple. Apply
Theorem~\ref{thm:cm-transfer}.

For different parameters, use instead the relation
$d_s(a,b)\leq2^{-U(h)}$. Net coverage again gives totality; the specified
parameter pairs now belong to this relation by hypothesis. The same
atomic bound and the same transfer theorem apply.
Every relation test and carrier search in this proof is finite and
rational.
\end{proof}

Fix an integer $\ell$ with $A=\max(1,A_\varphi)\leq2^\ell$ and define
\begin{equation}
 p_k=k+\ell+2,\qquad n_k=U(p_k)+1,\qquad
 v_\varphi(\bar a,k)=E_\varphi(\bar a;n_k,p_k).
 \label{eq:cm-canonical-value}
\end{equation}
Then $v_\varphi$ is primitive recursive. Indeed the syntax is fixed,
each atomic procedure and net generator is fixed, and every quantifier
is a bounded fold over a primitive-recursively generated list.
Nested quantifiers can make this computation extremely large; no
polynomial or elementary complexity bound is claimed.

\begin{corollary}[A real-valued interpretation with an iPRA Cauchy proof]
\label{cor:cm-cauchy}
For all dense parameters and all $k,l$,
\begin{equation}
 |v_\varphi(\bar a,k)-v_\varphi(\bar a,l)|
 \leq \tfrac12\bigl(2^{-k}+2^{-l}\bigr).
 \label{eq:cm-cauchy}
\end{equation}
For each fixed formula, the contracts
\eqref{eq:cm-net}--\eqref{eq:cm-atom-contract} therefore yield an iPRA
derivation of the regular-name condition for the fixed numerical
function $v_\varphi$.
\end{corollary}
\begin{proof}
In \eqref{eq:cm-cross-mesh} set $h=\min(p_k,p_l)$.
Monotonicity of $U$ makes its mesh hypotheses hold. Since
$2^{-h}=\max(2^{-p_k},2^{-p_l})\leq2^{-p_k}+2^{-p_l}$, the right side
is at most
$2A(2^{-p_k}+2^{-p_l})
\leq(2^{-k}+2^{-l})/2$.
This is a rational calculation for arbitrary numerical $k,l$.
No limit is taken inside the object theory.
\end{proof}

Define $\varphi^\CM(\bar a)$ externally to be the real represented by
this regular name. The comparison theorem gives two useful consequences:
\begin{equation}
\begin{aligned}
 |E_\varphi(\bar a;n,p)-\varphi^\CM(\bar a)|
   &\leq A_\varphi(2^{-p}+2^{-h})
       &&(n\geq U(h)),\\
 |\varphi^\CM(\bar a)-\varphi^\CM(\bar b)|
   &\leq A_\varphi2^{-h}
       &&(d_\infty(\bar a,\bar b)\leq2^{-U(h)}).
\end{aligned}
\label{eq:cm-semantic-bounds}
\end{equation}
For the first inequality, compare the given evaluation with
\eqref{eq:cm-canonical-value} at successively finer precisions using
\eqref{eq:cm-cross-mesh}, with $h$ fixed. The term $2^{-q}$ tends to
zero and all remaining terms are fixed.
The second follows similarly from the different-parameter estimate.
These are constructive properties of the represented real values.
They are semantic adequacy statements, separate from the finite
arithmetical theorem \eqref{eq:cm-cross-mesh}.

Thus $\mu_\varphi(k)=U(k+\ell)$ is a modulus on the dense presentation.
The completion-extension theorem, Theorem~\ref{thm:an-extension},
gives the safe completed-parameter modulus
\begin{equation}
 \nu_\varphi(k)=\mu_\varphi(k+2)+2=U(k+\ell+2)+2.
 \label{eq:cm-completed-modulus}
\end{equation}
The margin matters because a closed-threshold modulus on dense points
need not extend with its original indices.
For a tuple $\bar x$ of regular names, whose $j$th entries are
$\bar x_j$, an explicit output schedule is
\begin{equation}
 j_k=U(k+\ell+3)+2,\qquad p=k+\ell+2,\qquad n=U(p)+1,
 \qquad E_\varphi(\bar x_{j_k};n,p).
 \label{eq:cm-named-parameters}
\end{equation}
The distance from $\bar x_{j_k}$ to its completed tuple is at most
$2^{-j_k}$, so \eqref{eq:cm-completed-modulus} bounds the parameter
error by $2^{-(k+1)}$. The first inequality of
\eqref{eq:cm-semantic-bounds}, with $h=p$, bounds the evaluation
error by $2A_\varphi2^{-p}\leq2^{-(k+1)}$.
Thus the returned rational has total error at most $2^{-k}$.
This finite schedule is primitive recursive after fixing the input
name routines.
For arbitrary external names this is a finite-query construction
relative to the names; it is not a new first-order quantification over
all such names.

\subsection{The extension to distances given by regular names}
\label{subsec:cm-named-distance}

The exact rational pseudometric was a simplifying presentation choice.
The function sort in Section~\ref{sec:functionspaces} has a
primitive-recursive \emph{approximation} to its supremum distance;
the preceding theorem must therefore be extended before it is applied
to that sort. The required extension preserves a decidable comparison
relation by reserving an explicit margin.

Let $\Delta_s(a,b,r)$ be a symmetric rational distance sampler,
with values differing from the represented distance by at most $2^{-r}$.
Symmetry can be imposed by averaging the samplers at $(a,b)$ and $(b,a)$.
Its arithmetic coherence theorem is
\begin{equation}
 |\Delta_s(a,b,r)-\Delta_s(a,b,t)|\leq2^{-r}+2^{-t}.
 \label{eq:cm-distance-coherence}
\end{equation}
Define the first-order upper-bound formula
\begin{equation}
 \operatorname{Upper}_s(a,b,\rho)\quad\equiv\quad
 \forall r\ 
    \Delta_s(a,b,r)\leq \rho+2^{-r}.
 \label{eq:cm-upper}
\end{equation}
This formula is not decidable and will not be used as a scan condition.
Its semantic meaning is exactly $d_s(a,b)\leq\rho$.
Supply the net contract
$\forall a\ \exists j<|N_s(n)|\
\operatorname{Upper}_s(a,N_s(n)_j,2^{-n})$
and the atomic stability contract \eqref{eq:cm-atom-contract} with
the metric premise replaced by the conjunction of these
$\operatorname{Upper}$ formulas at radius $2^{-\mu_P(h)}$.
These are explicit object-theory obligations for the presentation.

\begin{lemma}[A decidable relation for named distances]
\label{lem:cm-named-relation}
Fix $h$, put $u=U(h)$ and $\delta=2^{-u}$, and take nets of indices
$n,m\geq u+2$. On their augmented finite carriers define
\begin{equation}
 aR_sb\quad\Longleftrightarrow\quad
 \Delta_s(a,b,u+3)\leq \tfrac58\delta.
 \label{eq:cm-slack-relation}
\end{equation}
Then $R_s$ is total both ways, identical free parameters are related,
and every related pair satisfies
$\operatorname{Upper}_s(a,b,\tfrac34\delta)$.
Consequently the cross-mesh estimate
\eqref{eq:cm-cross-mesh} holds for named distances under the stronger
mesh assumptions $n,m\geq U(h)+2$.
\end{lemma}
\begin{proof}
The other net covers a specified carrier point within
$2^{-m}\leq\delta/4$. Instantiate its $\operatorname{Upper}$
certificate at precision $u+3$:
the sampled distance to the covering center is at most
$\delta/4+\delta/8=3\delta/8$, and therefore passes
\eqref{eq:cm-slack-relation}. This proves each coverage direction;
symmetry handles orientation. A diagonal pair has sampled distance
at most $\delta/8$.

For a pair passing the test, coherence at an arbitrary precision $r$
gives
\[
 \Delta_s(a,b,r)
 \leq \Delta_s(a,b,u+3)+2^{-(u+3)}+2^{-r}
 \leq \tfrac34\delta+2^{-r}.
\]
Generalizing over $r$ proves the $\operatorname{Upper}$ assertion.
Since $\tfrac34\delta\leq2^{-\mu_P(h)}$, the supplied atomic stability
contract gives the atomic distortion used in
Theorem~\ref{thm:cm-transfer}. The relation test itself is rational
and primitive recursive. Although an all-precision upper bound is
derived, no induction formula contains that unbounded quantifier:
$r$ was arbitrary in one arithmetic inequality.
\end{proof}

In this version replace \eqref{eq:cm-canonical-value} by
\begin{equation}
 p_k=k+\ell+2,\qquad
 n_k=U(p_k)+2,\qquad
 v_\varphi(\bar a,k)=E_\varphi(\bar a;n_k,p_k).
 \label{eq:cm-named-schedule}
\end{equation}
The Cauchy calculation is unchanged. To compare different parameters,
a distance at most $2^{-(U(h)+1)}=\delta/2$ suffices:
its sampled distance is at most $5\delta/8$, so those parameter pairs
also pass the relation test. The resulting dense-parameter modulus is
$U(k+\ell)+1$. The safe modulus after completion is consequently
\begin{equation}
 \nu_\varphi^{\mathrm{named}}(k)=U(k+\ell+2)+3.
 \label{eq:cm-completed-named-modulus}
\end{equation}
For free parameter names, replace the schedule
\eqref{eq:cm-named-parameters} by
$j_k=U(k+\ell+3)+3$, $p=k+\ell+2$, and $n=U(p)+2$.
The same two errors, each at most $2^{-(k+1)}$, prove output
accuracy $2^{-k}$.
The first inequality of \eqref{eq:cm-semantic-bounds} remains valid
with $n\geq U(h)+2$.

The quantifier-adequacy proof below uses only this version of the
cross-mesh estimate, its uniform parameter modulus, and finite net
coverage. In the named-distance case take $n=U(p)+2$ throughout
that proof: its net radius is at most $2^{-(U(p)+2)}$, which is
smaller than the parameter radius just specified. Thus the constants
in its displayed value and witness estimates are unchanged.
In particular, the rational supremum-distance sampler of
Lemma~\ref{lem:function-distance} and the envelope net satisfy exactly
the numerical interfaces required to use a compact function class as
a quantified sort. No exact comparison of two real distances is used.

\subsection{Why the interpretation really has infimum and supremum}
\label{subsec:cm-quantifier-adequacy}

The construction above did not assume that every bounded subset of the
reals has a supremum. We now prove that its numerical quantifier
operations have exactly the intended meanings.

\begin{theorem}[Constructive quantifier adequacy]
\label{thm:cm-adequacy}
The values just constructed respect every connective in
\eqref{eq:cm-syntax} and satisfy
\begin{equation}
 (\sup_{y:s}\psi)^\CM(\bar a)
    =\sup_{b\in\widehat D_s}\psi^\CM(\bar a,b),
 \qquad
 (\inf_{y:s}\psi)^\CM(\bar a)
    =\inf_{b\in\widehat D_s}\psi^\CM(\bar a,b).
 \label{eq:cm-adequacy}
\end{equation}
Here the displayed suprema and infima exist because the proof constructs
them. In addition, a primitive-recursive bounded scan supplies a dense
point whose value is within $2^{-k}$ of the requested extremum at
precision $k$.
\end{theorem}
\begin{proof}
For a connective, use one common sufficiently fine mesh and atomic
precision for its immediate subformulas. Their finite evaluations obey
the connective equation exactly. The uniform estimates
\eqref{eq:cm-semantic-bounds}, together with the connective's Lipschitz
bound, pass that equation to the represented values. It is immaterial
that the independently defined canonical names might use different
precisions: \eqref{eq:cm-cross-mesh} compares them.

Consider a supremum and write $S$ for its constructed value. At mesh
$n=U(p)+1$, let
\[
 M_{n,p}=\max_{b\in K_s(n,\bar a)}
                 E_\psi(\bar a,b;n,p),
\]
where all subformula quantifiers use the fixed carriers of this
evaluation. The extra-support version of
\eqref{eq:cm-semantic-bounds} yields, uniformly in these $b$,
\begin{equation}
 |E_\psi(\bar a,b;n,p)-\psi^\CM(\bar a,b)|
       \leq 2A_\psi2^{-p}.
 \label{eq:cm-child-error}
\end{equation}
Given any dense $c$, choose $b$ in the finite net with
$d_s(c,b)\leq2^{-n}\leq2^{-U(p)}$.
The second inequality of \eqref{eq:cm-semantic-bounds} and then
\eqref{eq:cm-child-error} give
\[
 \psi^\CM(\bar a,c)
 \leq\psi^\CM(\bar a,b)+A_\psi2^{-p}
 \leq M_{n,p}+3A_\psi2^{-p}.
\]
Taking arbitrarily fine $p$ proves $\psi^\CM(\bar a,c)\leq S$.
Uniform continuity extends this upper-bound property to every completed
point $c$.

Conversely, suppose $z$ bounds $\psi^\CM(\bar a,b)$ above at every
dense point. For every member of the finite net,
\eqref{eq:cm-child-error} gives
$E_\psi(\bar a,b;n,p)\leq z+2A_\psi2^{-p}$.
Taking the finite maximum and then the represented limit gives
$S\leq z$. Hence $S$ is the least upper bound. This proves existence
and the supremum equation without selecting an infinite sequence of
maximizers. The infimum proof reverses all inequalities.

Finally choose by a finite scan a $b_*$ attaining $M_{n,p}$.
The first estimate in \eqref{eq:cm-semantic-bounds}, applied to the
supremum formula, bounds $|S-M_{n,p}|$ by $2A_\psi2^{-p}$.
Together with \eqref{eq:cm-child-error},
\[
 S-\psi^\CM(\bar a,b_*)\leq4A_\psi2^{-p}.
\]
Choose $p=k+\ell+2$ with $A_\psi\leq2^\ell$. The right side is at most
$2^{-k}$. The maximizing index, its point code, and all precision
choices are primitive recursive. The minimizing case is identical.
\end{proof}

There need be no coherent choice of the approximate witnesses as $k$
varies. The theorem constructs a real value and finite witnesses for
each requested accuracy. Constructing an \emph{exact} optimizing point
is a further problem and may need a uniqueness or separation modulus.
This distinction is exactly what prevents the argument from
reintroducing an unproved branch-selection principle.

\subsection{A worked sentence with genuine quantifier alternation}

Take the sort $[0,1]$, rational dense points, dyadic grids, and the
atomic predicate $P(x,y)=|x-y|$. For this atom,
\[
 \bigl||x-y|-|x'-y'|\bigr|
 \leq |x-x'|+|y-y'|
 \leq2d_\infty((x,y),(x',y')),
\]
so $\mu_P(h)=h+1$ is valid. Consider
\begin{equation}
 \Phi=\inf_x\sup_y |x-y|.
 \label{eq:cm-radius-sentence}
\end{equation}
For a fixed $x$, the triangle inequality on the line shows
$|x-y|\leq\max(x,1-x)$ for $0\leq y\leq1$, while $y=0,1$
give equality. Also $\max(x,1-x)\geq1/2$ because one of $x,1-x$
is at least their average; at $x=1/2$ the maximum is $1/2$.
Consequently $\Phi=1/2$.

Every dyadic grid containing $0,1,1/2$ already returns this exact value
under the finite evaluator. More complicated alternating expressions
need not stabilize exactly, but the same cross-mesh estimate applies.
The logical construction handles $\inf\sup\inf\sup$ without allocating
an additional precision loss for each quantifier: the common net
correspondence handles all of them simultaneously. The price of more
quantifiers is the number of bounded evaluations, not a deterioration
of the proven error bound.

The sorts need not consist of scalar points. In particular, using
the named-distance extension in
Section~\ref{subsec:cm-named-distance}, the compact
class of bounded $L$-Lipschitz functions constructed elsewhere in this
report can be a sort. Evaluation $P(f,x)=f(x)$, after affine rescaling
to $[0,1]$, is uniformly continuous because
\[
 |f(x)-g(y)|\leq\|f-g\|_\infty+L\|x-y\|_\infty.
\]
The supplied finite nets consist of actual Lipschitz functions, so a
formula may quantify over that function class and over points, using
the same theorem. This is a controlled interpretation of apparently
higher-order analytic statements. It is restricted to the specified
compact equicontinuous class; it does not quantify over all continuous
functions or over all primitive-recursive programs.

\subsection{A witness form of the Tarski--Vaught argument}

Global compactness can be replaced by more local data when only a
specified family of formulas is of interest. Let $A$ be a nonempty
subdomain of a structure $M$, with all atomic interpretations obtained
by restriction. Fix a finite family of formulas closed under
subformulas. For every subformula $\inf_y\psi(\bar x,y)$, suppose a
specified finite nonempty list $W_{\psi,\bar a,k}\subseteq A$ satisfies
\begin{equation}
 \forall b\in M\ \exists w\in W_{\psi,\bar a,k}\quad
 \psi^M(\bar a,w)\leq\psi^M(\bar a,b)+2^{-k}.
 \label{eq:cm-witness-test}
\end{equation}
For a supremum use the reverse inequality. The list generator and
the proof of its contract are input data. In an arithmetic application
they must have numerical realizations and iPRA proofs; an informal
existence assertion is not a generator.

\begin{proposition}[Supplied witnesses yield elementary observations]
\label{prop:cm-tarski-vaught}
For the specified subformula family, the values computed over $A$
and over $M$ agree at every tuple from $A$, provided those values
are constructed. With a uniform witness tolerance $\tau$ instead of
arbitrarily small $2^{-k}$, the error for the core basis is at most
$q(\varphi)\tau$, where $q$ is defined by
\begin{equation}
\begin{aligned}
 q(P)=q(c)&=0,\qquad q(1-\varphi)=q(\varphi),\\
 q(\min/\max(\varphi,\psi))&=\max(q(\varphi),q(\psi)),\\
 q(\inf/\sup_y\varphi)&=q(\varphi)+1.
\end{aligned}
 \label{eq:cm-rank-budget}
\end{equation}
\end{proposition}
\begin{proof}
The atomic case is equality, and the connective steps follow from
nonexpansiveness. Suppose the subformula values differ by at most
$e$ on tuples from $A$, and consider an infimum. Use only the two
constructed values
\[
 L_A=\inf_{w\in A}\psi^A(\bar a,w),\qquad
 L_M=\inf_{b\in M}\psi^M(\bar a,b).
\]
For every $w\in A$,
$L_M\leq\psi^M(\bar a,w)\leq\psi^A(\bar a,w)+e$.
Therefore $L_M-e$ is a lower bound for all the values defining
$L_A$, and $L_M\leq L_A+e$.
For each $b\in M$, the witness contract supplies $w\in A$ with
\[
 L_A\leq\psi^A(\bar a,w)
      \leq\psi^M(\bar a,w)+e
      \leq\psi^M(\bar a,b)+\tau+e.
\]
Consequently $L_A-\tau-e$ is a lower bound for all the values
defining $L_M$, so $L_A\leq L_M+\tau+e$.
These two inequalities give the asserted error $e+\tau$.
No auxiliary infimum of $\psi^M$ over $A$ is assumed to exist.
The supremum argument is dual.
Induction gives \eqref{eq:cm-rank-budget}.
When the witness contract is available at every precision, let the
explicit tolerance tend to zero in the represented real values.
\end{proof}

This is weaker input than a metric net if a particular formula has an
efficient approximate minimizer even when the entire domain is not
totally bounded. It is also a more precise statement of the unresolved
obligation: one must produce and verify the lists in
\eqref{eq:cm-witness-test}. The theorem does not produce elementary
submodels from arbitrary existence claims, and its conclusion does
not supply a general first-order satisfaction decision procedure.

\begin{proposition}[Observable quotients]
\label{prop:cm-observable-quotient}
Suppose each sort has a projection $\pi_s:M_s\to K_s$ with a
specified section $j_s:K_s\to M_s$, so $\pi_sj_s$ is the identity.
Suppose every atomic recipe factors exactly through $\pi$:
\[
 P^M(\bar a)=P^K(\pi\bar a).
\]
Then every formula in the chosen language has
$\varphi^M(\bar a)=\varphi^K(\pi\bar a)$ whenever the values are
constructed. If $K$ has the compact presentation above, its finite
evaluator supplies these values even when $M$ itself is not compact.
\end{proposition}
\begin{proof}
For an atomic formula this is the factorization hypothesis.
Connectives preserve equality. For a quantifier the two sets of
subformula values coincide: projection sends every $M$-argument to a
$K$-argument, and the supplied section lifts every $K$-argument back.
Therefore their constructed suprema and infima coincide. Equivalently,
the graph of $\pi$ is a total correspondence of atomic error zero.
This proves the result without choosing any representative not already
provided by the section.
\end{proof}

For a constructive noncompact example, let $M$ be the closed unit ball
of $c_0$: its elements are real sequences $(x_j)$ with $|x_j|\leq1$,
supplied with a modulus $b_x(k)$ such that
$j\geq b_x(k)$ implies $|x_j|\leq2^{-k}$.
The supremum distance is a represented real, not an assumed supremum
of an arbitrary bounded set. To approximate it to precision $k$,
take the finite prefix ending before
$B=\max(1,b_x(k+2),b_y(k+2))$. At each of its coordinates,
approximate $x_j,y_j$ separately to error $2^{-(k+3)}$ and take
the absolute difference of the two rational samples. These are
nonnegative approximants to $|x_j-y_j|$ with error at most
$2^{-(k+2)}$. Their finite maximum
approximates the full distance to error at most $2^{-(k+1)}$:
outside the prefix every difference is at most $2^{-(k+1)}$,
and inside it the specified sampling bound applies. Rescaling by
$1/2$ gives diameter at most one if required.

Let $K=[-1,1]^d$. Projection takes the first $d$ coordinates; the
section pads with zeros and supplies the constant tail bound $d$.
The domain $M$ is not totally bounded: the coordinate unit vectors
are pairwise at supremum distance one, so a ball of radius less than
$1/2$ contains at most one of them. Nevertheless every formula whose
atoms depend only on the first $d$ coordinates reduces exactly to the
compact cube $K$. Rational coordinate formulas therefore have the
finite arithmetic interpretation above. Adding the full
supremum-distance predicate would invalidate the factorization
hypothesis. The scope of the observable language is part of the theorem.

\subsection{Exact formalization boundary and the remaining Lean work}
\label{subsec:cm-lean}

The supplied file \texttt{ContinuousSyntax.lean} defines arity-indexed
syntax \path{QuantitativeIPRA.Continuous.Formula}, a finite numerical
\texttt{Model}, environment extension, a finite evaluator, and
\texttt{Correspondence}. Its theorem
\path{QuantitativeIPRA.Continuous.formula_transfer} proves the
unchanged-error result for the full core basis, including arbitrary
alternation of finite infima and suprema. It uses the checked
\texttt{min\_correspondence} and \texttt{max\_correspondence}
lemmas in \texttt{FiniteExtrema.lean}. All carriers are nonempty ranges
$0,\ldots,\mathit{last}$. A rational score is represented by a natural
numerator over a common positive denominator. The predicate
\[
 \operatorname{Close}(e,a,b)\quad\equiv\quad
       a\leq b+e\ \wedge\ b\leq a+e
\]
avoids an incorrect use of truncated natural subtraction for absolute
error. Truncated complement agrees with $B-a$ on $0\leq a\leq B$.

These are checked Lean propositions about finite semantic structures,
\emph{not} terms of the repository's
$\texttt{Prf PRAAx}$ object theory. The full compact-presentation
compiler and its semantic adequacy theorem have the human proofs above
and the following explicit formalization path:
\begin{enumerate}
\item Introduce a record \texttt{CompactPresentation} containing the
      dense-code validity predicate, rational pseudometric, net lengths
      and entries, and object-theory derivations of
      \eqref{eq:cm-net}. Introduce \texttt{AtomicRecipe} with fixed
      $\PR$ terms for its evaluator and modulus, together with
      derivations of \eqref{eq:cm-atom-contract}. No field is a
      universal evaluator for arbitrary programs.
\item Compile each fixed formula to numerical $\PR$ terms using
      composition and the bounded folds already used for
      \texttt{ApproximateWitness}. For the core basis use one positive
      denominator, normalize rational constants, and compile complements
      as bounded subtraction.
\item Prove object-theory finite correspondence transport. The existing
      \texttt{FiniteNet}, \texttt{NetTransfer}, and
      \texttt{ApproximateWitness} modules supply net elimination and
      extremal-index infrastructure. Reproduce the four rational
      inequalities in the proof of
      Theorem~\ref{thm:cm-transfer}, then assemble one derivation per
      fixed syntax tree by the usual substitution and cut operations.
\item Define the finite relation
      $d_s(a,b)\leq2^{-U(h)}$ and prove its two coverage directions
      from \eqref{eq:cm-net}. Instantiate the object correspondence
      theorem to derive \eqref{eq:cm-cross-mesh}, and then
      \eqref{eq:cm-cauchy}. These are arithmetic endpoints with every
      mesh, precision, parameter and denominator explicit.
\item In a separate semantic namespace, define the real value by the
      regular Cauchy name. Prove \eqref{eq:cm-semantic-bounds} and
      Theorem~\ref{thm:cm-adequacy}, using the completion-extension
      theorem. A Mathlib realization may use completed metric spaces
      here, but that realization is an adequacy theorem and must not
      be relabeled an iPRA derivation.
\end{enumerate}

The fixed-formula proof transformation is enough for the proposed
applications. A uniform compiler on numerical syntax codes would
require its own primitive-recursive size bounds, proof-code encoding,
and correctness theorem. The checked Lean induction above is useful
evidence for that future implementation; it does not by itself supply
those additional object-theory results.

Finally, zero-valued truth remains different from finite observation.
The algorithm returns intervals of decreasing width containing
$\varphi^\CM(\bar a)$. It can certify a strict inequality when the
interval separates from its threshold. It need not decide the
statement $\varphi^\CM(\bar a)=0$. Neither equality of arbitrary real
names nor the complete theory of an arbitrary presented structure
has become decidable.
